\documentclass[10pt,a4paper]{amsart}
\usepackage[margin=19mm]{geometry}
\usepackage{amsmath,amssymb,mathtools}
\usepackage[scr=rsfs,cal=euler]{mathalfa}
\usepackage{xfrac}
\usepackage[unicode,hidelinks,bookmarks=false]{hyperref}
\newcommand{\ie}{i.e.\ }
\newcommand{\cf}{cf.\ }
\newcommand{\resp}{respectively\ }
\newcommand{\Subsection}{\subsection}
\providecommand{\nobreakdash}{\nobreak\hbox{-}}
\setlength{\emergencystretch}{2em}
\multlinegap0pt
\usepackage{amssymb}
\usepackage[all]{xy}
\usepackage{xcolor}
\usepackage[shortlabels]{enumitem}
\newtheorem{lem}{Lemma}
\newtheorem{overpic}{Overpic}
\title{Stein traces and characterizing slopes}
\author{Roger Casals, John Etnyre,, Marc Kegel}
\date{Source: arXiv:2111.00265v3 (CC BY 4.0)}
\begin{document}
\maketitle

\noindent\emph{Source and licence.} Stein traces and characterizing slopes. Version arXiv:2111.00265v3 (CC BY 4.0), distributed by arXiv under the Creative Commons Attribution 4.0 International licence, http://creativecommons.org/licenses/by/4.0/, as recorded in the arXiv licence metadata for this exact version and shown on the abstract page https://arxiv.org/abs/2111.00265v3. Full licence notice: \texttt{licenses/arxiv-2111.00265-CC-BY-4.0.txt}.\par\smallskip

\noindent\textit{Editorial scope.} Counted unit: the article's lem secondannulus (source0.tex lines 704--706). The article's complete original proof of it is delivered with it (source0.tex lines 709--745, one \texttt{proof} environment of 37 lines). No mathematics is written, repaired or re-derived: the statement and the proof are the article's own lines, in the article's own notation, and no retained line is substituted or re-flowed.  No further source-local context is needed: every cross-reference of every retained line resolves inside the two delivered spans.  Nothing was omitted from any retained span, so the package carries no omission record.  No retained line contains a citation, so the wrapper carries no bibliography and no bibliography entry.  No retained line contains a figure environment, an image-inclusion command or drawing code, so no image is delivered and the package declares no asset dependency.  Editorial formatting only, in addition: an AMS \texttt{amsart} preamble that restates the article's own declarations and macros used by the retained lines, namely the environments \texttt{lem}, \texttt{overpic} on separate counters, together with the standard packages those lines need.  The retained environments print in this document as Lemma 1, Overpic 1 in order of appearance, whereas the article prints the same items with the numbering of its own full document.  The complete native source of the article as distributed in the arXiv e-print for this exact version is bundled unchanged as \texttt{sources/117877/source0.tex}.
\begin{lem}\label{secondannulus}
 The Legendrian knots $K_1$ and $K_2$, considered as Legendrian knots in the surgered contact manifold $L_{A,\gamma}(-1)$, are Legendrian isotopic.
\end{lem}

\begin{proof}
 We represent our Legendrian knots in surgery diagrams for $(M,\xi)$, drawn as front projections, and the Legendrian arc $\gamma$ will be attached to $K$ and $K'$ at cusps in the front diagram. Figure~\ref{fig:part1} (top) shows the Legendrian knot $L_{A,\gamma}$ and the two Legendrian knots, $K_1$ and $K_2$, that bound the pre-Lagrangian sub-annulus $A'$.  We perform a contact $(-1)$-surgery along $L_{A,\gamma}$. After performing a handle slide of $K_2$ over $L_{A,\gamma}$, we obtain Figure~\ref{fig:part1} (bottom). Note that the Legendrian knot $L_{A,\gamma}$ is made out of four Legendrian arcs: $\widetilde{K'}$ and $\widetilde{K}$, where $\widetilde{K}$ denotes $K$ with a small neighborhood of the $\partial \gamma$ removed from it (and similarly for $\widetilde{K'}$), and two copies of $\gamma$, that we denote $\gamma'$ and $\gamma''$.
 
	\begin{figure}[htbp]{\small
			\begin{overpic}%[grid,tics=10] 
				{figs/atwistpart1.pdf}
				\put(109, 143){$\gamma$}
				\put(11, 172){\color{red} $K_2$}
				\put(11, 117){\color{blue} $K_1$}
				\put(158, 110){$(-1)$}
				\put(70, 164){$L_{A,\gamma}$}
				\put(109, 41){$\gamma$}
				\put(70, 65){$L_{A,\gamma}$}
				\put(158, 8){$(-1)$}
		\end{overpic}}
		\caption{The top diagram is $L_{A,\gamma}$ together with the Legendrian knots $K_1$ and $K_2$. The bottom figure is the result of Legendrian surgery on $L_{A,\gamma}$  and sliding $K_2$ over $L_{A,\gamma}$.}
		\label{fig:part1}
	\end{figure}


	
The Legendrian knot $K_2$, after the slide, and the Legendrian arc $\widetilde{K'}$ are joined by a cusp, as depicted in the upper left of the front diagram in Figure~\ref{fig:part1} (bottom). Since $K_2$ and $K'$ are parallel, there is a Legendrian isotopy ending in the upper diagram of Figure~\ref{fig:part2}. 
	\begin{figure}[htbp]{\small
			\begin{overpic}%[grid,tics=10] 
				{figs/atwistpart2.pdf}
				\put(109, 144){$\gamma$}
				\put(109, 41){$\gamma$}
				\put(70, 169){$L_{A,\gamma}$} 
				\put(70, 62){$L_{A,\gamma}$}
				\put(160, 8){$(-1)$}
				\put(160, 110){$(-1)$}
		\end{overpic}}
		\caption{Steps in the Legendrian isotopy from $K_2$ to $K_1$, as used in the proof for Lemma \ref{secondannulus}.}
		\label{fig:part2}
	\end{figure}
Now, the Legendrian curves $\gamma'$ and $\gamma''$ are Legendrian push-offs of $\gamma$, and connected with a cusp as in the upper diagram in Figure~\ref{fig:part2}. Thus, there is a Legendrian isotopy to the bottom diagram in Figure~\ref{fig:part2}, which is isotopic to $K_1$, as required. 
\end{proof}

\end{document}
